\section{Introduction}
Astrophotography, as a hobby, sits at an unusual intersection between art and measurement. It is often approached as a craft guided by experience, software tools, and the knowledge of other imagers. This empirical method can be quite effective: many excellent images are produced by following established recipes or adopting the techniques of skilled processors. However, recipes alone do not always explain why a method works, when it fails, or what assumptions it makes about the data.\\[\baselineskip]
This becomes especially important as processing tools grow more sophisticated. Modern, AI-based deconvolution, denoising, and sharpening tools can produce substantial improvements, but their inner workings are difficult to inspect. As a result, we often apply powerful transformations to our astronomical data without a solid understanding of what they do, what assumptions they rely on, or how they alter the underlying signal. This intuitive approach may be acceptable when the aim is to produce aesthetically pleasing images; for quantitative work or scientific analysis, it is simply insufficient.
\subsection{Scope and assumptions}
This manual focuses primarily on image processing and therefore generally omits topics related to image acquisition, such as telescope choice, mounts, guiding, filters, and sensors, except where they directly affect a processing method (e.g. DRIZZLE). The reader is thus assumed to be familiar with the basic terminology and practical aspects of astrophotography.\\[\baselineskip]
As for the mathematical side, we assume comfort with basic Calculus, Linear Algebra, and probability distributions. Technical observations and clarifications are placed in specially marked boxes in order to preserve the main flow of the exposition. These boxes are used for subtleties, rigorous definitions, and optional remarks that may deepen understanding but are not essential on a first reading.